\section{Introduction}

Thousands of approved drugs act on human biology in ways that only become
apparent once the drugs are in widespread use, producing both adverse
reactions and, occasionally, unexpected therapeutic benefits that motivate
drug repurposing. Large pharmacovigilance resources such as SIDER catalog
many of these associations by aggregating side effects and indications at
scale \cite{kuhn2016sider}, while DrugBank and the Drug Repurposing Hub
annotate the molecular targets and clinical uses of thousands of
compounds \cite{knox2024drugbank,corsello2017repurposing}. Molecular
profiling resources further characterize how drugs perturb biology: the EPA
ToxCast/Tox21 program assays hundreds of compounds across more than a
thousand biological endpoints \cite{richard2016toxcast}, and the
Connectivity Map and its L1000 successor profile drug-induced gene-expression
signatures \cite{lamb2006cmap,subramanian2017lincs}. Collectively these
resources raise the prospect of reasoning systematically about which disease
pathways a drug influences and why.

In parallel, phenome-wide genetic studies are now producing gene-level maps
of disease risk at unprecedented breadth. The UK Biobank cohort has enabled
genome-wide association studies of thousands of traits
\cite{bycroft2018ukbiobank}, and gene-based methods such as PrediXcan,
S-PrediXcan and S-MultiXcan integrate GWAS summary statistics with
tissue-specific eQTL models to map genetic associations from variants to
genes \cite{gamazon2015predixcan,barbeira2018spredixcan,barbeira2019smultixcan}.
PhenomeXcan compiles S-PrediXcan/S-MultiXcan results for thousands of
UK Biobank phenotypes into a queryable gene-by-phenotype resource
\cite{pividori2020phenomexcan}. On the tissue-regulatory side, the Genotype-Tissue
Expression project provides an atlas of genetic regulatory effects
across 49 tissues that anchors these transcriptome-wide analyses \cite{gtex2020atlas}.

Most prior computational methods for discovering unexpected drug effects
have exploited one of these data modalities at a time. A dominant family of
approaches uses the connectivity principle -- that drugs effective for a
disease should induce gene-expression changes that oppose the disease's
molecular signature -- either directly from expression data
\cite{lamb2006cmap,sirota2011drugrepo} or after converting GWAS to
imputed transcriptomes for selected disorders
\cite{so2017neuropsychiatric,reay2020pes,grotzinger2023tsem}. A second family
uses drug-side-effect similarity to infer shared targets or repurposing
candidates \cite{campillos2008sideeffect}. A third family treats drug-disease
prediction as matrix completion on the incomplete drug-phenotype matrix.
These approaches have successfully generated repositioning candidates, but
they are typically predictive black boxes: they do not explain how the
molecular effects of a drug propagate through disease driver genes to
produce the observed phenotypic effects, and therefore do not yield the
mechanistic hypotheses that experimentalists and clinicians need.

We propose to close this gap by learning a single shared network that
connects drug molecular effects directly to the genes that drive disease
risk. Our method, Draphnet (Drug and Phenotype Network), is a supervised
logistic affinity-regression model that takes three matrices as input: a
drug-by-endpoint matrix $D$ compiled from ToxCast \cite{richard2016toxcast},
a gene-by-phenotype matrix $P$ compiled from PhenomeXcan
\cite{pividori2020phenomexcan,barbeira2019smultixcan}, and a drug-by-phenotype
matrix $Y$ compiled from SIDER \cite{kuhn2016sider}. Draphnet learns an
interaction matrix $W$ such that $DWP^{\top}$ approximates the logit of $Y$,
imposes sparsity through $L_1$ regularization, and handles the high
dimensionality of $D$ and $P$ through a SoftImpute decomposition of $D$
\cite{mazumder2010softimpute} and a singular value decomposition of $P$, a
strategy adapted from affinity regression in transcriptional
biology \cite{osmanbeyoglu2014affinity}. The same network is learned once
across thousands of drug-phenotype pairs, enforcing the prior that the
biological mechanisms connecting endpoints to disease genes are largely
shared across contexts. Phenotype names from UK Biobank and SIDER are
harmonized through UMLS concept unique identifiers \cite{bodenreider2004umls}.

Our contributions are threefold. First, we demonstrate that molecular and
genetic similarity alone already carry strong signal for drug-disease
biology: drugs with similar ToxCast endpoint profiles share significantly
more phenotypes than expected ($p=1\text{e-}22$; linear model adjusted for
shared-endpoint count $p=1.7\text{e-}43$), and diseases with more similar
PhenomeXcan profiles are perturbed by more similar sets of drugs
($p=4\text{e-}81$). Second, we show that Draphnet's learned model generalizes
to held-out drugs and significantly outperforms a nearest-neighbor baseline
(rank-sum $p=3\text{e-}8$). Third, and most importantly, by projecting each
drug through the learned interaction matrix into a 10,027-gene disease-genome
space, we recover mechanistically interpretable target-level enrichments --
28 of 132 DrugBank targets shared by at least three drugs have disease-gene
associations significant at adjusted $p<0.01$ -- and surface unexpected
drug-gene associations including PPM1M for HTR2C-modulating neuroleptics
and CETP for fenofibrate and other PPAR$\alpha$-targeting drugs, the latter
consistent with experimental evidence that PERM1-family factors
co-regulate PPAR$\alpha$ target genes \cite{huang2022perm1}.

\section{Related work}

Several lines of prior work motivate the components of Draphnet but do not,
to our knowledge, integrate them. Drug-repurposing methods based on the
connectivity principle \cite{lamb2006cmap,subramanian2017lincs,sirota2011drugrepo}
have repeatedly shown that shared gene-expression responses can link drugs
and diseases, but they operate on gene-expression signatures rather than on
disease-driver genetics, and they do not learn an interaction network.
Extensions of this principle to genetically informed transcriptomes using
S-PrediXcan have been explored for neuropsychiatric disorders
\cite{so2017neuropsychiatric,reay2020pes} and, more recently, through
transcriptome-wide structural equation modeling across 13 psychiatric
disorders \cite{grotzinger2023tsem}, but these analyses are restricted to a
narrow disease family and report candidate drugs rather than drug-gene
networks.

A complementary line uses drug-phenotype similarity to infer drug-target
relationships. Campillos et al.\ showed that side-effect similarity can
predict shared molecular targets \cite{campillos2008sideeffect}, and this
family of methods has since been extended with matrix-completion and
network-based approaches. These methods predict either missing drug-disease
annotations or shared targets, but, again, they do not expose the underlying
drug-endpoint-to-disease-gene wiring.

Finally, methodological work on polypharmacology has made clear that most
clinically effective drugs act on multiple proteins and that effective
biological models must allow for this complexity
\cite{hopkins2008network}. This observation is directly relevant to the
way we interpret Draphnet: drug targets such as CHRM1, which is annotated as
a target for 14 mechanistically diverse compounds, are unlikely to produce
homogeneous phenome-effect signatures, and our method is designed to
tolerate such heterogeneity through target-specific rather than global
enrichment tests. On the statistical side, the empirical permutation null
and Benjamini-Yekutieli correction \cite{benjamini2001fdrdependency,benjamini1995fdr}
allow us to control false discoveries across $\sim$10,000 correlated
gene-level tests per drug, and graph analyses on disease-gene overlap are
performed with NetworkX \cite{hagberg2008networkx}.
